This thesis developed an automated, LLM-powered static analysis tool to detect SQL antipatterns in Java source code that uses jOOQ, and subsequently conducted a large-scale empirical analysis of their estimated prevalence.

The study began with mining \num{602} open-source jOOQ projects. To establish a ground truth, 61 representative projects were manually annotated, resulting in \num{1562} antipattern occurrences across 19 SQL antipatterns. Following statistical viability filtering, 12 antipatterns were excluded, leaving seven for all subsequent stages of the study. This filtered dataset enabled a systematic evaluation of four advanced LLMs across multiple prompting strategies. The evaluation revealed that complex reasoning prompts, such as Chain-of-Thought and Tree-of-Thought, generally failed to improve detection performance. Instead, these strategies drastically increased computational costs, extended runtimes, and in some cases degraded model stability. Consequently, the final detection tool was engineered using Claude Opus 4.5 without reasoning, with a baseline Zero-Shot Prompting strategy, achieving a high F1-score of \num{0.88} in a multi-label multi-occurrence localisation task.

Deploying the finalised tool across the entire \num{602}-project corpus flagged \num{15931} occurrences of seven distinct antipatterns. The empirical data demonstrated frequent co-occurrence among detected antipatterns; notably, the “Implicit Columns” and “ID Required” antipatterns were almost ubiquitous, occurring in nearly \qty{90}{\percent} of the analysed repositories. Furthermore, the analysis revealed that the vast majority of the detected query antipattern occurrences were associated with a few specific jOOQ convenience methods, such as \texttt{selectFrom} and \texttt{Field.like}.

To ensure repeatability and facilitate future research, we open-sourced all artefacts produced as parts of the thesis. The finalised tool is available in the following GitHub repository: \href{https://github.com/kristoisberg/jooq-antipattern-detector}{https://github.com/kristoisberg/jooq-antipattern-detector}. All other artefacts, including experimental scripts and their results, evaluated prompts, and source code for this document, are available in the following GitHub repository: \href{https://github.com/kristoisberg/masters-thesis}{https://github.com/kristoisberg/masters-thesis}.

Ultimately, this study highlights the potential of LLMs to address the parsing limitations of dynamic, DSL-based database frameworks. By detailing the convenience methods most commonly observed alongside detected antipatterns, the research provides actionable insights to help developers avoid common database pitfalls and recognise recurring design risks in their codebases.
